% ---------------------------------------------------------------------
%  Chapitre 4 — Le contexte des trades (version aperçu, rédigée)
%  NB : chiffres de l'ancienne base (38 035), à régénérer sans l'oracle.
% ---------------------------------------------------------------------
\chapter{Le contexte des trades~: comparer et caractériser}
\label{chap:contexte}

Le chapitre précédent a jugé chaque stratégie isolément, contre le
hasard. Celui-ci change de question. Les stratégies diffèrent-elles
entre elles~? Le succès dépend-il du contexte, du régime de marché, du
secteur, de la situation d'entrée~? Les conclusions de ce chapitre sont
comparatives et descriptives~: on n'y valide aucune stratégie, le
verdict appartenant au chapitre~\ref{chap:juger}.

\begin{limites}
La stratégie oracle (témoin, chapitre~\ref{chap:juger}) est exclue de ce
chapitre~: son avantage artificiel (\edgevar{} moyen de l'ordre de
$+26\,\%$) écraserait toute comparaison. Elle n'a de sens que comme
témoin de validation, non comme stratégie à comparer aux autres. Les
analyses ci-dessous portent donc sur les dix stratégies réelles.
\end{limites}

% =====================================================================
\section{Les stratégies diffèrent-elles entre elles ?}
\label{sec:b1-etape2}

Pour comparer les stratégies entre elles, on analyse la variance de
l'\edgevar{}. Si les moyennes des dix stratégies étaient égales, la
variation entre stratégies ne dépasserait pas la variation interne à
chacune. La statistique \(F\) rapporte ces deux variations~:
\[
  F = \frac{\mathrm{SCE}_{\text{inter}}/(k-1)}
           {\mathrm{SCE}_{\text{intra}}/(n-k)} \sim F_{k-1,\,n-k}.
\]
La décomposition de la variance et le test \(F\) suivent le cadre
classique de l'analyse de la variance~\cite{scheffe1959}. Les
comparaisons deux à deux qui suivent le rejet global s'appuient sur le
test de Tukey, fondé sur la loi du range studentisé, qui corrige de
lui-même la multiplicité des paires. Une analyse à deux facteurs,
croisant la stratégie et le régime de marché avec un terme
d'interaction, teste enfin si l'effet de la stratégie dépend du régime.

\begin{conditions}
Indépendance des observations, hypothèse mise à l'épreuve au
chapitre~\ref{chap:juger}~: les p-values de ce chapitre, calculées au
niveau des trades, s'entendent sous cette réserve, les \(F\) colossaux
restant interprétables comme ordres de grandeur~; normalité des résidus
assurée par le théorème central limite~; homoscédasticité vérifiée par
les tests de Bartlett et de Levene, avec bascule possible vers l'ANOVA
de Welch en cas de variances inégales.
\end{conditions}

L'analyse à un facteur rejette massivement l'égalité des moyennes, avec
\(F\approx 47\) et une p-value de l'ordre de $10^{-85}$. Le test de
Tukey identifie dix-huit paires significatives sur quarante-cinq,
\texttt{rsi\_strict} se détachant du lot. L'interaction entre stratégie
et régime est très significative, de l'ordre de $5\times10^{-69}$~: la
valeur d'une stratégie dépend du régime dans lequel elle opère.

% =====================================================================
\section{Gagner dépend-il du contexte ?}
\label{sec:b1-etape3}

On se demande ensuite si le simple fait de gagner dépend du contexte. La
variable \texttt{win}, qui vaut un pour un trade gagnant, est croisée
avec le régime, le secteur et la stratégie dans des tableaux de
contingence, sur lesquels on applique le test du khi-deux
d'indépendance~:
\[
  \chi^2 = \sum_{i,j}\frac{(N_{ij}-E_{ij})^2}{E_{ij}},
  \qquad E_{ij}=\frac{N_{i\cdot}N_{\cdot j}}{N},
  \qquad \chi^2 \sim \chi^2_{(r-1)(c-1)}.
\]

\begin{conditions}
Effectifs théoriques \(E_{ij}\ge 5\) (règle de Cochran)~; sur tableaux
restreints ou à faibles effectifs, recours possible à la correction de
Yates, au test exact de Fisher ou au G-test~; indépendance des
observations, même réserve qu'à la section précédente.
\end{conditions}

\begin{table}[H]
\centering
\caption{Dépendance de \texttt{win} au contexte (khi-deux d'indépendance).}
\label{tab:b1-etape3}
\small
\begin{tabular}{@{}lrrc@{}}
\toprule
Croisement & $\chi^2$ & ddl & Verdict \\
\midrule
\texttt{win} $\times$ \texttt{regime\_entry} & $112{,}3$   & 2  & lien significatif \\
\texttt{win} $\times$ \texttt{sector}        & $16{,}6$    & 10 & indépendance ($p=0{,}08$) \\
\texttt{win} $\times$ \texttt{strategy}      & $6\,381{,}7$& 9  & lien très fort \\
\bottomrule
\end{tabular}
\end{table}

Le taux de gain dépend fortement de la stratégie, faiblement du régime,
et pas du secteur.

\begin{limites}
À très grand \(n\), le khi-deux devient significatif même pour des liens
négligeables~; la significativité ne dit rien de la force du lien, et
ces contrastes se lisent comme des descriptions, non comme des
validations.
\end{limites}

% =====================================================================
\section{Le contexte d'entrée en carte --- ACP}
\label{sec:b1-etape4}

Pour cartographier les situations d'entrée, on réalise une ACP
normée~\cite{saporta2006,lebart2006} sur les quatre variables
quantitatives de contexte, la volatilité, le niveau du RSI, la distance
à la moyenne mobile et la durée de détention, le nombre d'axes retenus
étant justifié par les critères de Kaiser et de Cattell. Les variables
qualitatives, régime, secteur et issue du trade, sont projetées en
illustratives.

\begin{conditions}
Variables quantitatives~; standardisation obligatoire, les unités étant
hétérogènes~; l'ACP ne capture que des relations linéaires.
\end{conditions}

Les deux premiers axes résument environ deux tiers de la variance, le
premier étant porté par le RSI et la distance à la moyenne mobile,
autrement dit par l'état de survente du titre à l'entrée, le second par
la durée de détention et la volatilité. La question intéressante est de
savoir si les trades gagnants occupent une région particulière de ce
plan. La réponse est négative~: les positions moyennes des gagnants et
des perdants sont presque confondues, avec un écart standardisé
inférieur à $0{,}2$. Le contexte d'entrée, à lui seul, ne sépare pas les
gagnants des perdants, résultat instructif qui rejoint le verdict du
chapitre~\ref{chap:juger}.

\begin{figure}[H]
  \centering
  \includegraphics[width=0.9\linewidth]{etape4_acp.png}
  \caption{ACP du contexte d'entrée~: plan factoriel des trades, gagnants
  et perdants presque confondus.}
  \label{fig:acp-contexte}
\end{figure}

% =====================================================================
\section{AFC et ACM}
\label{sec:b1-etape56}

L'analyse factorielle des correspondances croise ensuite la stratégie
avec la tranche de rendement obtenue, perte, gain modéré ou gain fort.
Le premier axe, qui résume les trois quarts de l'inertie, oppose les
stratégies à dominante perdante, comme \texttt{hs\_classic} ou
\texttt{sr\_breakout}, aux stratégies dont les trades finissent plus
souvent en gain, \texttt{rsi\_classic} et surtout \texttt{rsi\_strict},
cette dernière se projetant nettement du côté des gains forts.

L'analyse des correspondances multiples, menée sur la stratégie, le
régime, le secteur et l'issue du trade, ne résume que faiblement
l'information, comme il est habituel avec de nombreuses modalités, mais
confirme la même géographie~: les modalités de secteur se massent au
centre du plan, signe qu'elles ne structurent rien, tandis que les
stratégies et le régime s'écartent, \texttt{rsi\_strict} voisinant avec
le régime baissier, dont elle exploite les excès.

\begin{figure}[H]
  \centering
  \includegraphics[width=0.85\linewidth]{etape5_afc.png}
  \caption{AFC stratégie $\times$ tranche de rendement~: le premier axe
  oppose profils perdants et profils gagnants.}
  \label{fig:afc}
\end{figure}

% =====================================================================
\section{Synthèse du chapitre}
\label{sec:contexte-synthese}

Les stratégies diffèrent donc fortement entre elles et leur valeur
dépend du régime de marché, tandis que le secteur ne joue aucun rôle
décelable. Mais aucun de ces contrastes ne remet en cause le verdict du
chapitre~\ref{chap:juger}~: différer de son voisin n'est pas battre le
hasard. Ces cartes du contexte préparent en revanche la perspective d'un
avantage conditionnel, esquissée au chapitre~\ref{chap:protocole}.
